\documentclass[../../main.tex]{subfiles}% 注意这里的文件路径不能用 ./main.tex ,否则用latexmk编译子文件会报错
\graphicspath{{\subfix{./image/}}{\subfix{../../image/}}} % 指定图片引用路径,后续可以直接使用图片文件名
% 如果图片存放在上级目录,则使用 ../../image/ 作为备用路径

% 例如:
% \begin{figure}[H]
% \centering
% \includegraphics[width=0.3\textwidth]{图.png}
% \caption{}
% \label{figure:图}
% \end{figure}
% 注意:上述\label{}一定要放在\caption{}之后,否则引用图片序号会只会显示??.

\begin{document}

\section{多元函数的连续性和微分}

我们的极限采用聚点定义，即只需要沿着有定义的地方趋近即可.

\subsection{基本问题}

\begin{example}
设\(f(x,y)=\begin{cases}\dfrac{x^2 y}{x^4 + y^2},&x^2 + y^2 \neq 0\\0,&x^2 + y^2 = 0\end{cases}\)，证明\(f\)沿着每条射线\(\begin{cases}x = t\cos\alpha,&t > 0,\alpha \in [0,2\pi)\\y = t\sin\alpha,&t > 0,\alpha \in [0,2\pi)\end{cases}\)趋于\((0,0)\)时都趋于 0，但是\(f\)在\((0,0)\)不连续.
\end{example}
\begin{note}
本结果表明，使用极坐标求二重极限不一定正确.实际上,我们用极坐标求二重根极限时,都是固定$\alpha$,再令$t\to 0$求极限.因此得到的只是沿着每个过原点的射线(与$x$轴的夹角为$\alpha$)趋于$(0,0)$的极限，比如还可以沿$y=kx^2$这条曲线趋于$(0,0)$.
\end{note}
\begin{proof}
一方面，
\begin{align*}
\lim_{t \to 0^+} f(t\cos\alpha, t\sin\alpha)&=\lim_{t \to 0^+} \frac{t^3 \cos\alpha \sin\alpha}{t^4 \cos^4 \alpha + t^2 \sin^2 \alpha} = \lim_{t \to 0^+} \frac{t \cos\alpha \sin\alpha}{t^2 \cos^4 \alpha + \sin^2 \alpha} = 0,
\end{align*}
另外一方面
\[
\lim_{x \to 0^+} f(x, kx^2)=\lim_{x \to 0^+} \frac{kx^4}{x^4 + k^2 x^4}=\frac{k}{1 + k^2}.
\]
故\(f\)在\((0,0)\)不连续.矛盾!
\end{proof}
\begin{remark}
实际上,使用极坐标变换求极限时,只需要在$t\to 0$的时候让$\alpha$也发生变化(不再固定$\alpha$),再求极限才能得到正确的极限值,但这样反而不方便求极限.

那么什么时候固定$\alpha$后求出来的极限就是原函数的极限呢？实际上只需要极限关于\(\alpha \in [0,2\pi)\)一致，因为你直接考察定义\(\varepsilon - \delta\)语言即可. 实际做题中可以体现为
\[
\lim_{t \to 0^+} \sup_{\alpha \in [0,2\pi)} |f(t\cos\alpha, t\sin\alpha) - A| = 0
\]
更直白的，你需要得到形如
\[
|f(t\cos\alpha, t\sin\alpha) - A| \leqslant g(t)
\]
的不等式且\(\lim_{t \to 0^+} g(t) = 0\).
\end{remark}

\begin{proposition}\label{proposition:二重极限极坐标换元的条件}
设二元函数 \( f(x, y) \) 在点 \( (a, b) \) 的某个去心邻域内有定义,若对$\forall \alpha \in [0, 2\pi)$,都有
\begin{align*}
\underset{t \rightarrow 0^+}{\lim}f\left( a+t\cos \alpha ,b+t\sin \alpha \right) =A\in \mathbb{R},
\end{align*}
并且
\begin{align*}
\lim_{t \to 0^+} \sup_{\alpha \in [0,2\pi)} |f(t\cos\alpha, t\sin\alpha) - A| = 0,
\end{align*}
或者
存在函数 \(g(t)\) 满足对$\forall \alpha \in [0, 2\pi)$,都有
\[
|f(a + t\cos\alpha, b + t\sin\alpha) - A| \leqslant g(t),\quad \lim_{t \to 0^+} g(t) = 0,
\]
则
\[
\lim_{(x, y) \to (a, b)} f(x, y) = A.
\]
\end{proposition}
\begin{proof}
显然条件$\lim_{t \to 0^+} \sup_{\alpha \in [0,2\pi)} |f(t\cos\alpha, t\sin\alpha) - A| = 0$和条件
存在函数 \(g(t)\) 满足对$\forall \alpha \in [0, 2\pi)$,都有
\[
|f(a + t\cos\alpha, b + t\sin\alpha) - A| \leqslant g(t),\quad \lim_{t \to 0^+} g(t) = 0
\]
等价.
由$\lim\limits_{t\rightarrow 0^+}g\left( t \right) =0$知，$\forall \varepsilon >0$，存在$\delta >0$，使得
\begin{align*}
|g(t)| < \varepsilon,\ \forall t \in (0, \delta).
\end{align*}
对$\forall (x, y)$，满足$0 < \sqrt{(x - a)^2 + (y - b)^2} < \delta$，令
\begin{align*}
t = \sqrt{(x - a)^2 + (y - b)^2},\ \alpha = \arctan\frac{y - b}{x - a},
\end{align*}
则$x = a + t\cos\alpha$，$y = b + t\sin\alpha$。于是
\begin{align*}
|f(x, y) - A| = |f(a + t\cos\alpha, b + t\sin\alpha) - A| \leqslant g(t) < \varepsilon,\ \forall (x, y) \in B\left( (a, b), \delta \right).
\end{align*}
\end{proof}

\begin{example}
计算
\[
\lim_{(x,y) \to (0,0)} \frac{x^3 + y^3}{x + y}
\]
\end{example}
\begin{proof}
考虑
\[
|f(t\cos\alpha, t\sin\alpha)| = t^2 \left| \frac{\cos^3 \alpha + \sin^3 \alpha}{\cos\alpha + \sin\alpha} \right| = t^2 |\cos^2 \alpha + \sin^2 \alpha - \cos\alpha \sin\alpha| \leqslant 2t^2,
\]
于是
\[
0 \leqslant \lim_{t \to 0^+} |f(t\cos\alpha, t\sin\alpha)| \leqslant 2 \lim_{t \to 0^+} t^2 = 0
\]
故我们得到了
\[
\lim_{(x,y) \to (0,0)} \frac{x^3 + y^3}{x + y} = 0.
\]
\end{proof}

\begin{example}
设\(f\)在\((0,0)\)连续且满足
\[
\lim_{(x,y) \to (0,0)} \frac{f(x,y) - xy}{x^2 + y^2} = a > 0
\]
求\(a\)的范围使得\(f\)在\((0,0)\)一定取到极值. 再求\(a\)的范围使得\(f\)在\((0,0)\)一定取不到极值.
\end{example}
\begin{note}
注意到\(\lim_{(x,y) \to (0,0)} \frac{xy}{x^2 + y^2}\)不存在，但是
\[
\left| \frac{xy}{x^2 + y^2} \right| \leqslant \frac{1}{2}
\]
即猜测\(\frac{1}{2}\)是\(a\)的分界点.
\end{note}
\begin{proof}
由条件容易得到
\begin{align*}
&\quad \quad \,\,\underset{\left( x,y \right) \rightarrow \left( 0,0 \right)}{\lim}\frac{f\left( x,y \right) -xy}{x^2+y^2}=a
\\
&\Longrightarrow \underset{\left( x,y \right) \rightarrow \left( 0,0 \right)}{\lim}f\left( x,y \right) =a\cdot \underset{\left( x,y \right) \rightarrow \left( 0,0 \right)}{\lim}\left( x^2+y^2 \right) +\underset{\left( x,y \right) \rightarrow \left( 0,0 \right)}{\lim}xy=0
\\
&\Longrightarrow f(0,0)=\lim_{(x,y)\rightarrow (0,0)} f(x,y)=0
\end{align*}
以及
\[
f(x,y) = (a + g(x,y))(x^2 + y^2) + xy, \quad \lim_{(x,y) \to (0,0)} g(x,y) = 0.
\]
当\(a > \frac{1}{2}\)，当\((x,y)\)足够靠近\(0\)使得\(g(x,y) > \frac{1}{2} - a\). 此时我们有
\[
f(x,y) = (a + g(x,y))(x^2 + y^2) + xy > \frac{1}{2}(x^2 + y^2) + xy = \frac{1}{2}(x + y)^2 \geqslant 0
\]
故\(a > \frac{1}{2}\)时,$f$在$(0,0)$处取得极小值.

当\(0 < a < \frac{1}{2}\)，当\((x,y)\)足够靠近\(0\)使得\(-a < g(x,y) < \frac{1 - 2a}{4}\)，则此时当\(x > 0, y > 0\)有\(f(x,y) > 0\). 但是
\[
f(x,y) = (a + g(x,y))(x^2 + y^2) + xy < \frac{1 + 2a}{4}(x^2 + y^2) + xy
\]
又\(\frac{1 + 2a}{4}(x^2 + y^2) + xy\)在\(y = -x\)上有
\[
\frac{1 + 2a}{4}2x^2 - x^2 = \frac{2a - 1}{2}x^2 < 0,
\]
故此时\((0,0)\)不是$f$的极值点.

当\(a = \frac{1}{2}\)问题无法判断，这是因为考虑\(f(x,y) = xy + \frac{1}{2}(x^2 + y^2) + x^2(x^2 + y^2)\)，则
\[
f(x,y) = \frac{1}{2}(x + y)^2 + x^2(x^2 + y^2) > 0,
\]
即\((0,0)\)是极值. 但是考虑\(f(x,y) = \frac{1}{2}(x^2 + y^2) + xy - x(x^2 + y^2)\)，就有
\[
f(x,-x) = x^2 - x^2 - 2x^3 = -2x^3 < 0, x > 0,
\]
\[
f(x,x) = x^2 + x^2 - 2x^3 = 2x^2 - 2x^3 > 0, 0 < x < 1.
\]
即\((0,0)\)不是极值.
\end{proof}

\begin{theorem}\label{theorem:定理雅克比行列式为0的充要条件}
设 \(u=f(x,y),\ v=g(x,y)\) 在区域 \(D\subset\mathbb R^2\) 上有连续偏导数。对任意点 \(P_0=(x_0,y_0)\in D\)，下列结论成立：
\begin{enumerate}[(1)]
\item 若 \(J(P_0)=\dfrac{\partial(u,v)}{\partial(x,y)}(P_0)=0\)，则存在 \(P_0\) 的一个邻域 \(U\subset D\)，使得
\begin{enumerate}[(i)]
\item 若 \(\nabla v(P_0)\ne(0,0)\)，则存在\(U\)上的一元连续可微函数 \(F\)，满足在 \(U\) 上 \(u=F(v)\)；
\item 若 \(\nabla u(P_0)\ne(0,0)\)，则存在\(U\)上的一元连续可微函数 \(G\)，满足在 \(U\) 上 \(v=G(u)\)；
\item 若 \(\nabla u(P_0)=\nabla v(P_0)=(0,0)\)，则 \(u,v\) 在 \(U\) 上均为常数，从而可选取任意一元连续可微函数使得$u,v$之间有函数关系.
\end{enumerate}
\item 反之，若存在 \(P_0\) 的某个邻域 \(U\) 和一元连续可微函数 \(F\) 或 \(G\)，使得在 \(U\) 上 \(u=F(v)\) 或 \(v=G(u)\)，则 \(J(P_0)=0\)。
\end{enumerate}
\end{theorem}
\begin{proof}
\begin{enumerate}[(1)]
\item  设 \(J(P_0)=0\)。若 \(\nabla u(P_0)=\nabla v(P_0)=(0,0)\)，则由偏导数的连续性，在 \(P_0\) 的某邻域内 \(u,v\) 均为常数，结论显然成立。否则，至少有一个梯度非零。

先考虑 \(\nabla v(P_0)\ne(0,0)\) 的情形。不妨设 \(v_y(P_0)\ne0\),若 \(v_y(P_0)=0\)，则必有 \(v_x(P_0)\ne0\)，此时交换 \(x,y\) 的角色即可。由\hyperref[theorem:隐函数可微性定理]{隐函数可微性定理}，存在 \(P_0\) 的邻域 \(U\) 以及连续可微函数 \(y=\psi(x,v)\)，使得在 \(U\) 上恒有
\begin{align*}
v=g(x,\psi(x,v)).
\end{align*}
定义
\begin{align*}
\Phi(x,v)=f(x,\psi(x,v)),
\end{align*}
则在 \(U\) 上 \(u=\Phi(x,v)\)。由于 \(J=0\)，计算雅可比行列式：
\begin{align*}
0=J=
\begin{vmatrix}
u_x & u_y\\
v_x & v_y
\end{vmatrix}
=
\begin{vmatrix}
\Phi_x+\Phi_v v_x & \Phi_v v_y\\
v_x & v_y
\end{vmatrix}
=\Phi_x v_y+\Phi_v v_xv_y-\Phi_v v_yv_x=\Phi_x v_y.
\end{align*}
因为 \(v_y\ne0\)，故 \(\Phi_x=0\)。这表明 \(\Phi\) 在 \(U\) 上不依赖于 \(x\)，即存在一元连续可微函数 \(F\) 使得 \(u=\Phi(v)=F(v)\)。

若 \(\nabla v(P_0)=0\) 但 \(\nabla u(P_0)\ne0\)，则类似地，不妨设 \(u_y(P_0)\ne0\)。由\hyperref[theorem:隐函数可微性定理]{隐函数可微性定理}得 \(y=\phi(x,u)\)，令
\begin{align*}
\Psi(x,u)=g(x,\phi(x,u)),
\end{align*}
则在 \(U\) 上 \(v=\Psi(x,u)\)。同样利用 \(J=0\)，有
\begin{align*}
0=J=
\begin{vmatrix}
u_x & u_y\\
v_x & v_y
\end{vmatrix}
=
\begin{vmatrix}
u_x & u_y\\
\Psi_x+\Psi_u u_x & \Psi_u u_y
\end{vmatrix}
= u_x\Psi_u u_y - u_y(\Psi_x+\Psi_u u_x) = -u_y\Psi_x,
\end{align*}
因 \(u_y\ne0\)，得 \(\Psi_x=0\)，故 \(\Psi\) 与 \(x\) 无关，即存在一元连续可微函数 \(G\) 使 \(v=\Psi(u)=G(u)\)。
\item 若在 \(U\) 上有 \(u=F(v)\)，则对 \(U\) 内任意点求偏导，得
\begin{align*}
u_x=F'(v)v_x,\qquad u_y=F'(v)v_y,
\end{align*}
因此
\begin{align*}
J=u_xv_y-u_yv_x=F'(v)v_xv_y-F'(v)v_yv_x=0.
\end{align*}
同理，若 \(v=G(u)\)，则
\begin{align*}
v_x=G'(u)u_x,\qquad v_y=G'(u)u_y,
\end{align*}
从而
\begin{align*}
J=u_xv_y-u_yv_x=u_xG'(u)u_y-u_yG'(u)u_x=0.
\end{align*}
\end{enumerate}
\end{proof}

\begin{example}
设 \( u = f(x, y), v = g(x, y) \) 在区域 \( \Omega \subset \mathbb{R}^2 \) 上有连续偏导数。证明：
如果在 \( \Omega \) 上的 Jacobi 矩阵的秩恒等于 \( 1 \)，则对 \( \Omega \) 内任一点 \( (x_0, y_0) \)，存在 \( (x_0, y_0) \) 的邻域 \( U \) 和连续可微的函数 \( F(u, v) \)，且 \( (F'_u)^2 + (F'_v)^2 \neq 0 \)，使
\[
F(f(x, y), g(x, y)) = 0
\]
在 \( U \) 上恒成立。
\end{example}
\begin{proof}
由题设，Jacobi矩阵的秩恒为 \( 1 \)，故其行列式
\begin{align*}
J = \frac{\partial (u, v)}{\partial (x, y)} = 0,
\end{align*}
并且在每一点处至少有一个梯度非零。

任取 \( (x_0, y_0) \in \Omega \)。
若 \( \nabla v(x_0, y_0) \ne (0, 0) \)，则由\cref{theorem:定理雅克比行列式为0的充要条件}，在 \( (x_0, y_0) \) 的某个邻域 \( U \) 内存在一元连续可微函数 \( \varphi \)，使得
\begin{align*}
u = f(x, y) = \varphi(g(x, y)) = \varphi(v).
\end{align*}
令
\begin{align*}
F(u, v) = u - \varphi(v),
\end{align*}
则 \( F \) 在相应邻域内连续可微，并且对任意 \( (x, y) \in U \)，有
\begin{align*}
F(f(x, y), g(x, y)) = f(x, y) - \varphi(g(x, y)) = 0.
\end{align*}
同时
\begin{align*}
F'_u = 1, \qquad F'_v = -\varphi'(v),
\end{align*}
故
\begin{align*}
(F'_u)^2 + (F'_v)^2 = 1 + (\varphi'(v))^2 > 0.
\end{align*}
若 \( \nabla u(x_0, y_0) \ne (0, 0) \)，则由\cref{theorem:定理雅克比行列式为0的充要条件}可知存在一元连续可微函数 \( \psi \)，使得
\begin{align*}
v = g(x, y) = \psi(f(x, y)) = \psi(u)
\end{align*}
在邻域 \( U \) 内成立。令
\begin{align*}
F(u, v) = v - \psi(u),
\end{align*}
则
\begin{align*}
F(f(x, y), g(x, y)) = g(x, y) - \psi(f(x, y)) = 0,
\end{align*}
且
\begin{align*}
F'_u = -\psi'(u), \qquad F'_v = 1,
\end{align*}
从而
\begin{align*}
(F'_u)^2 + (F'_v)^2 = (\psi'(u))^2 + 1 > 0.
\end{align*}
综上，无论哪种情形，都存在所需的邻域 \( U \) 和函数 \( F \)，使得复合函数在 \( U \) 上恒为零.
\end{proof}

\begin{example}
设 \( x f'_x + y f'_y = 0 \)，证明 \( f \) 是 \( \frac{y}{x} \) 的函数.
\end{example}
\begin{proof}
注意到
\[
\begin{vmatrix}
f'_x & f'_y \\
-\frac{y}{x^2} & \frac{1}{x}
\end{vmatrix} = \frac{1}{x^2} \left( x f'_x + y f'_y \right) = 0, \frac{\partial \left( \frac{y}{x} \right)}{\partial x} = -\frac{y}{x^2}, \frac{\partial \left( \frac{y}{x} \right)}{\partial y} = \frac{1}{x}
\]
我们由\refthe{theorem:定理雅克比行列式为0的充要条件}知 \( f \) 是 \( \frac{y}{x} \) 的函数.
\end{proof}

\begin{example}
可微函数 $z = f(x, y)$ 仅是 $ax + by$（$ab \neq 0$）的函数的充分必要条件是
\begin{align*}
b \frac{\partial z}{\partial x} = a \frac{\partial z}{\partial y}.
\end{align*}
\end{example}
\begin{proof}
令 $u = z = f(x, y)$，$v = ax + by$。则 $\nabla v = (a, b)$，由 $ab \neq 0$ 知 $\nabla v \neq (0,0)$ 恒成立。计算Jacobi行列式：
\begin{align*}
J = \frac{\partial (u, v)}{\partial (x, y)}
= \begin{vmatrix}
z_x & z_y \\
a & b
\end{vmatrix}
= b z_x - a z_y.
\end{align*}
由\refthe{theorem:定理雅克比行列式为0的充要条件}知，存在一元可微函数 $F$ 使得 $z = F(ax+by)$的充要条件是
\begin{align*}
J = \frac{\partial (u, v)}{\partial (x, y)}
= \begin{vmatrix}
z_x & z_y \\
a & b
\end{vmatrix}
= b z_x - a z_y=0\iff b z_x = a z_y.
\end{align*}
\end{proof}

\begin{definition}
我们称 \( f : \mathbb{R}^2 \to \mathbb{R} \) 为\textbf{齐} \( \boldsymbol{n}( n \in \mathbb{N} )\) \textbf{次函数}，如果 \( f \) 满足
\[
f(tx, ty) = t^n f(x, y), \forall x, y \in \mathbb{R}, t > 0.
\]
\end{definition}

\begin{proposition}[齐次函数基本性质]\label{proposition:齐次函数基本性质}
若 \( f \in D^2(\mathbb{R}^2) \), 则 \( f \) 是齐 \( n \) 次函数的充要条件是
\[
x \frac{\partial f}{\partial x} + y \frac{\partial f}{\partial y} = n f.
\]
\end{proposition}
\begin{proof}
若 \( f \) 是齐 \( n \) 次函数, 则
\[
f(tx, ty) = t^n f(x, y), \forall x, y \in \mathbb{R}, t > 0.
\]
两边对 \( t \) 求导得
\[
x \frac{\partial f}{\partial x}(tx, ty) + y \frac{\partial f}{\partial y}(tx, ty) = n t^{n - 1} f(x, y),
\]
于是
\[
tx \frac{\partial f}{\partial x}(tx, ty) + ty \frac{\partial f}{\partial y}(tx, ty) = n t^n f(x, y) = n f(tx, ty),
\]
再令$\mathbf{x}=tx,\mathbf{y}=ty$,即证.

反过来若
\[
x \frac{\partial f}{\partial x} + y \frac{\partial f}{\partial y} = n f
\]
成立, 固定 \( x, y \in \mathbb{R} \) 并考虑 \( g(t) = f(tx, ty) \). 则有
\[
t g'(t) = tx \frac{\partial f}{\partial x}(tx, ty) + ty \frac{\partial f}{\partial y}(tx, ty) = n f(tx, ty) = n g(t).
\]
故解微分方程得 \( g(t) = C t^n \), 从而将$g(1) = f(x, y)$代入得$C=f(x,y),$于是
\[
g(t)=f(tx, ty) = t^n f(x, y).
\]
这就证明了
\[
f(tx, ty) = t^n f(x, y), \forall x, y \in \mathbb{R}, t > 0.
\]
\end{proof}

\begin{example}
设 \( \frac{\partial^2 u}{\partial x \partial y} + \frac{\partial u}{\partial y} = 0 \) 且 \( u(0, y) = y^2, u(x, 1) = \cos x \), 求 \( u \).
\end{example}
\begin{proof}
对 \( \frac{\partial^2 u}{\partial x \partial y} + \frac{\partial u}{\partial y} = 0 \) 两边关于 \( y \) 积分得 \( \frac{\partial u}{\partial x} + u = C(x) \). 由 \( u(x, 1) = \cos x \) 得
\[
\frac{\partial u}{\partial x}(x,1)=-\sin x\Rightarrow \frac{\partial u}{\partial x}\left( x,1 \right) +u\left( x,1 \right) =C(x)=\cos x-\sin x.
\]
又
\[
\frac{\partial (u e^x)}{\partial x} = e^x \left( \frac{\partial u}{\partial x} + u \right) = e^x (\cos x - \sin x) \Rightarrow u e^x = e^x \cos x + C_2(y),
\]
我们有
\[
u(x, y) = \cos x + C_2(y) e^{-x}.
\]
现在
\[
y^2 = u(0, y) = 1 + C_2(y) \Rightarrow C_2(y) = y^2 - 1.
\]
故
\[
u(x, y) = \cos x + (y^2 - 1) e^{-x}.
\]
\end{proof}

\begin{example}
设 \( l_1, l_2 \) 夹角为 \( \varphi \in (0, \pi) \) 且 \( f \) 连续可微, 证明
\[
\left| \frac{\partial f}{\partial x} \right|^2 + \left| \frac{\partial f}{\partial y} \right|^2 \leqslant \frac{2}{\sin^2 \varphi} \left[ \left| \frac{\partial f}{\partial l_1} \right|^2 + \left| \frac{\partial f}{\partial l_2} \right|^2 \right] .
\]
\end{example}
\begin{proof}
由可微时方向导数计算公式有
\[
\frac{\partial f}{\partial l_1} = \cos a \cdot \frac{\partial f}{\partial x} + \sin a \cdot \frac{\partial f}{\partial y},
\]
\[
\frac{\partial f}{\partial l_2} = \cos (a + \varphi) \cdot \frac{\partial f}{\partial x} + \sin (a + \varphi) \cdot \frac{\partial f}{\partial y},
\]
则
\begin{align*}
\begin{large}
\begin{pmatrix}
{\textstyle \frac{\partial f}{\partial l_1} }\\
{\textstyle \frac{\partial f}{\partial l_2}}
\end{pmatrix}
\end{large}
=
\begin{pmatrix}
\cos a & \sin a \\
\cos (a + \varphi) & \sin (a + \varphi)
\end{pmatrix}
\begin{large}
\begin{pmatrix}
{\textstyle \frac{\partial f}{\partial x}} \\
{\textstyle \frac{\partial f}{\partial y}}
\end{pmatrix}
\end{large}.
\end{align*}
于是
\[
\begin{aligned}
&\left| \frac{\partial f}{\partial l_1} \right|^2+\left| \frac{\partial f}{\partial l_2} \right|^2=\left( \frac{\partial f}{\partial l_1}\quad \frac{\partial f}{\partial l_2} \right) \left( \begin{array}{c}
\frac{\partial f}{\partial l_1}\\
\frac{\partial f}{\partial l_2}\\
\end{array} \right) 
\\
&=\left( \frac{\partial f}{\partial x}\quad \frac{\partial f}{\partial y} \right) \left( \begin{matrix}
\cos a&		\sin a\\
\cos\mathrm{(}a+\varphi )&		\sin\mathrm{(}a+\varphi )\\
\end{matrix} \right) ^T\left( \begin{matrix}
\cos a&		\sin a\\
\cos\mathrm{(}a+\varphi )&		\sin\mathrm{(}a+\varphi )\\
\end{matrix} \right)\begin{large}
\begin{pmatrix}
{\textstyle \frac{\partial f}{\partial x}} \\
{\textstyle \frac{\partial f}{\partial y}}
\end{pmatrix}
\end{large}
\\
&=\left( \frac{\partial f}{\partial x}\quad \frac{\partial f}{\partial y} \right) \left( \begin{matrix}
\cos ^2a+\cos ^2(a+\varphi )&		\sin\mathrm{(}a+\varphi )\cos\mathrm{(}a+\varphi )+\sin a\cos a\\
\sin\mathrm{(}a+\varphi )\cos\mathrm{(}a+\varphi )+\sin a\cos a&		\sin ^2a+\sin ^2(a+\varphi )\\
\end{matrix} \right) \begin{large}
\begin{pmatrix}
{\textstyle \frac{\partial f}{\partial x}} \\
{\textstyle \frac{\partial f}{\partial y}}
\end{pmatrix}
\end{large} .
\end{aligned}
\]
利用
\[
\begin{pmatrix}
\cos^2 a + \cos^2 (a + \varphi) & \sin (a + \varphi)\cos (a + \varphi) + \sin a \cos a \\
\sin (a + \varphi)\cos (a + \varphi) + \sin a \cos a & \sin^2 a + \sin^2 (a + \varphi)
\end{pmatrix}
\]
的特征值为 \( 1 \pm \cos \varphi \) 和\hyperref[Basis of Algebra-proposition:Rayleigh-quotient瑞丽商的基本性质]{Rayleigh quotient(瑞丽商)的基本性质}, 我们知道
\[
\left| \frac{\partial f}{\partial x} \right|^2 + \left| \frac{\partial f}{\partial y} \right|^2 \leqslant \frac{1}{1 - |\cos \varphi|} \left[ \left| \frac{\partial f}{\partial l_1} \right|^2 + \left| \frac{\partial f}{\partial l_2} \right|^2 \right] \leqslant \frac{2}{\sin^2 \varphi} \left[ \left| \frac{\partial f}{\partial l_1} \right|^2 + \left| \frac{\partial f}{\partial l_2} \right|^2 \right],
\]
即证.上式最后一个不等式是因为
\begin{align*}
&\quad \quad \,\, \frac{1}{1 - |\cos \varphi|} \geqslant \frac{2}{\sin^2 \varphi} \\
&\Longleftrightarrow 1 - |\cos \varphi| \leqslant \frac{\sin^2 \varphi}{2} \\
&\Longleftrightarrow 2 - 2|\cos \varphi| \leqslant 1 - \cos^2 \varphi \\
&\Longleftrightarrow \cos^2 \varphi - 2|\cos \varphi| + 1 \geqslant 0 \\
&\Longleftrightarrow (|\cos \varphi| - 1)^2 \geqslant 0.
\end{align*}
\end{proof}

\begin{example}
设 \( D \) 为单位圆盘, 考虑 \( f \in C^1(D) \cap C(\overline{D}) \) 且 \( |f| \leqslant 1 \), 证明: 
\begin{enumerate}[(1)]
\item 存在 \( D \)内部中的一个点 \( (x_0, y_0) \) 使得
\[
\left| \frac{\partial f}{\partial x}(x_0, y_0) \right|^2 + \left| \frac{\partial f}{\partial y}(x_0, y_0) \right|^2 \leqslant 16 .
\]

\item 存在 \( D \)内部中的一个点 \( (x_0, y_0) \) 使得\[
\left| \frac{\partial f}{\partial x}(x_0, y_0) \right|^2 + \left| \frac{\partial f}{\partial y}(x_0, y_0) \right|^2 \leqslant 1,
\]
并说明右端常数 \(1\) 是最佳的.
\end{enumerate}
\end{example}
\begin{note}
摄动想法, 考虑$g(x,y)=f(x,y)+\varepsilon (x^2+y^2)$,其中$\varepsilon$待定.此外很多同学疑惑构造函数咋来的,实际上这是完全没有必要的! 因为大家几乎都是记的.
\end{note}
\begin{proof}
\begin{enumerate}[(1)]
\item 考虑 \( g(x, y) = f(x, y) + 2(x^2 + y^2) \), 则由$|f|\leqslant 1$知\( g|_{\partial D} \geqslant 1, g(0, 0)=f(0,0) \leqslant 1 \). 故 \( g \) 最小值在 \( D \) 内取到.又由$g\in C(D)$,从而存在 \( D \) 中的一个 \( g \) 的最小值点 \( (x_0, y_0) \) 使得 \( \frac{\partial g}{\partial x}(x_0, y_0) = 0, \frac{\partial g}{\partial y}(x_0, y_0) = 0 \), 即
\[
\frac{\partial f}{\partial x}(x_0, y_0) = -4x_0, \frac{\partial f}{\partial y}(x_0, y_0) = -4y_0,
\]
这就得到了证明.

\item 将 \(f\) 换成 \(-f\) 不影响题设条件与待证结论，因此不妨设 \(f(0,0) \geqslant 0\).
考察函数
\begin{align*}
F(x,y) = f(x,y) - \sqrt{x^2+y^2}.
\end{align*}
在圆周 \(x^2+y^2=1\) 上，\(F(x,y) = f(x,y) - 1 \leqslant 0\)，而 \(F(0,0) = f(0,0) \geqslant 0\). 若 \(F\) 的最大值$F(x_1,y_1)$大于 \(0\)，则最大值点不可能在圆周上,否则
\begin{align*}
F(x_1,y_1)=f(x_1,y_1)-1\leqslant 0,
\end{align*}
矛盾!
若最大值等于 \(0\)，则原点就是一个最大值点. 因此，\(F\) 总能在圆盘内部的某点 \((x_0,y_0)\) 处取得最大值.
\begin{enumerate}[(i)]
\item 若 \((x_0,y_0) \neq (0,0)\)，由极值的必要条件知$F_x(x_0,y_0)=F_y(x_0,y_0)=0$,即
\begin{align*}
f_x(x_0,y_0) = \frac{x_0}{\sqrt{x_0^2+y_0^2}}, \quad f_y(x_0,y_0) = \frac{y_0}{\sqrt{x_0^2+y_0^2}},
\end{align*}
从而 \(f_x^2(x_0,y_0) + f_y^2(x_0,y_0) = 1\).

\item 若 \((x_0,y_0) = (0,0)\)，则对任意满足 \(u^2+v^2=1\) 的 \((u,v)\) 以及充分小的 \(t>0\)，由原点是 \(F\) 的最大值点可得
\begin{align*}
f(tu,tv)-t=F\left( tu,tv \right) \leqslant F\left( 0,0 \right) =f(0,0)\Longleftrightarrow \frac{f(tu,tv)-f(0,0)}{t}\leqslant 1.
\end{align*}
令 \(t \to 0^+\)，得到 \(f_x(0,0)u + f_y(0,0)v \leqslant 1\).

若 \((f_x(0,0), f_y(0,0)) \neq (0,0)\)，取
\begin{align*}
(u,v) = \frac{(f_x(0,0), f_y(0,0))}{\sqrt{f_x^2(0,0)+f_y^2(0,0)}},
\end{align*}
便有 \(\sqrt{f_x^2(0,0)+f_y^2(0,0)} \leqslant 1\); 若两个偏导数均为 \(0\)，结论显然成立. 因此总能找到圆盘内一点 \((x_0,y_0)\)，使得 \(f_x^2(x_0,y_0) + f_y^2(x_0,y_0) \leqslant 1\).
\end{enumerate}
最后取 \(f(x,y)=x\). 它满足 \(|f(x,y)| \leqslant 1\)，并且在整个圆盘上恒有 \(f_x^2(x,y) + f_y^2(x,y) = 1\). 因此右端常数不能小于 \(1\)，故 \(1\) 是最佳常数.
\end{enumerate}
\end{proof}

\begin{proposition}\label{proposition:极坐标变换偏导数相关结论}
设$f(x,y)$是$D$上的二元函数且偏导数都存在,令$x=r\cos \theta$,$y=r\sin \theta$,\( g(r, \theta) = f(r \cos \theta, r \sin \theta) \), 则
\begin{gather*}
r \frac{\partial g}{\partial r} = x \frac{\partial f}{\partial x} + y \frac{\partial f}{\partial y},
\\
\frac{\partial g}{\partial \theta} = -y \frac{\partial f}{\partial x} + x \frac{\partial f}{\partial y}.
\end{gather*}
\end{proposition}
\begin{proof}
直接求导得证.
\end{proof}

\begin{example}
设 \( \lim\limits_{r = \sqrt{x^2 + y^2} \to +\infty} \left( x \frac{\partial f}{\partial x} + y \frac{\partial f}{\partial y} \right) = a > 0 \), 证明 \( f \) 在 \( \mathbb{R}^2 \) 取得最小值.
\end{example}
\begin{remark}
本题关键是有一个隐藏条件:条件极限关于角度 \( \theta \) 的一致性.
\end{remark}
\begin{note}
积累想法设 \( g(r, \theta) = f(r \cos \theta, r \sin \theta) \), 则
\[
r \frac{\partial g}{\partial r} = x \frac{\partial f}{\partial x} + y \frac{\partial f}{\partial y}, \frac{\partial g}{\partial \theta} = -y \frac{\partial f}{\partial x} + x \frac{\partial f}{\partial y}.
\]
即联想\refpro{proposition:极坐标变换偏导数相关结论}.
\end{note}
\begin{proof}
注意到
\begin{align*}
\lim\limits_{r \to +\infty} r g'_r = \lim\limits_{r \to +\infty} r \frac{\partial g}{\partial r}= a > 0,
\end{align*}
于是存在 \( r_0 > 0 \) 使得 \( g'_r > 0, \forall r \geqslant r_0 \). 则此时
\[
g(r, \theta) \geqslant g(r_0, \theta), \forall r \geqslant r_0, \theta \in [0, 2\pi),
\]
故$g$的最小值在$D=\left\{ \left( r,\theta \right) :r\in \left[ 0,r_0 \right] ,\theta \in \left[ 0,2\pi \right) \right\}$取到,因此\( \min\limits_{\substack{r \in [0, r_0], \\ \theta \in [0, 2\pi]}} g(r, \theta) \) 为 \( g \) 最小值.
\end{proof}

\begin{example}
设 \( f \) 是 \( \mathbb{R}^2 \) 上的连续可微函数且 \( f(0, 1) = f(1, 0) \), 证明存在单位圆周上两个不同的点使得 \( y \frac{\partial f}{\partial x} = x \frac{\partial f}{\partial y} \).
\end{example}
\begin{note}
联想\refpro{proposition:极坐标变换偏导数相关结论}.
\end{note}
\begin{proof}
令$x=cos\theta ,y=\sin\theta $,考虑 \( g(\theta) = f(\cos \theta, \sin \theta) \), 注意到
\[
g(0) = g\left( \frac{\pi}{2} \right) = g(2\pi).
\]
由Rolle中值定理知,存在$\theta_1\ne \theta_2 \in [0,2\pi)$,记$x_i=\cos\theta_1,y_i=\sin\theta(i=1,2)$,使得
\begin{align*}
&\quad \quad \,\,g' \left( \theta _1 \right) =g' \left( \theta _2 \right) =0
\\
&\Longleftrightarrow \begin{cases}
-\sin \theta _1\frac{\partial f}{\partial x}\left( \cos \theta _1,\sin \theta _1 \right) +\cos \theta _1\frac{\partial f}{\partial y}\left( \cos \theta _1,\sin \theta _1 \right) =0,\\
-\sin \theta _2\frac{\partial f}{\partial x}\left( \cos \theta _2,\sin \theta _2 \right) +\cos \theta _2\frac{\partial f}{\partial y}\left( \cos \theta _2,\sin \theta _2 \right) =0.\\
\end{cases}
\\
&\Longleftrightarrow \begin{cases}
y_1\frac{\partial f}{\partial x}\left( x_1,y_1 \right) =x_1\frac{\partial f}{\partial y}\left( x_1,y_1 \right) ,\\
y_2\frac{\partial f}{\partial x}\left( x_2,y_2 \right) =x_2\frac{\partial f}{\partial y}\left( x_2,y_2 \right).\\
\end{cases}
\end{align*}
即单位圆周上有两个不同的点,使得 \( y \frac{\partial f}{\partial x} = x \frac{\partial f}{\partial y} \).
\end{proof}

\begin{example}
\begin{enumerate}
\item 设 \( f \in C^1(\mathbb{R}^2), f(0,0) = 0 \) 且 \( |\nabla f| \leqslant 1 \), 证明 \( |f(1,2)| \leqslant \sqrt{5} \).

\item 设 \( f \in C^1(\mathbb{R}^2), f(0,0) = 0 \) 且
\[
\left| \frac{\partial f}{\partial x} \right| \leqslant 2|x - y|, \left| \frac{\partial f}{\partial y} \right| \leqslant 2|x - y|,
\]
证明: \( |f(5,4)| \leqslant 1 \).
\end{enumerate}
\end{example}
\begin{remark}
\textbf{梯度}及其\textbf{模}定义为$\nabla f\triangleq \left( \frac{\partial f}{\partial x},\frac{\partial f}{\partial y} \right) ,\left| \nabla f \right|\triangleq \sqrt{\left( \frac{\partial f}{\partial x} \right) ^2+\left( \frac{\partial f}{\partial y} \right) ^2}.$
\end{remark}
\begin{remark}
第二问如果和上一问完全类似, 取积分路径为连接 \( (0,0),(5,4) \) 的线段,那么由第一型曲线积分和第二型曲面积分的联系和\hyperref[theorem:Cauchy不等式(离散版)]{Cauchy不等式(离散版本)}我们有
\[
\begin{aligned}
|f(5,4)|&=\left| \int_{(0,0)}^{(5,4)}{\left( \frac{\partial f}{\partial x}\mathrm{d}x+\frac{\partial f}{\partial y}\mathrm{d}y \right)} \right|=\left| \int_{(0,0)}^{(5,4)}{\left( \frac{\partial f}{\partial x}\cos \left( \widehat{\boldsymbol{t},x} \right) +\frac{\partial f}{\partial y}\sin \left( \widehat{\boldsymbol{t},x} \right) \right) \mathrm{d}s} \right|
\\
&\leqslant \int_{(0,0)}^{(5,4)}{|\nabla f|\mathrm{d}s}\leqslant \sqrt{8}\int_{(0,0)}^{(5,4)}{|x}-y|\mathrm{d}s
\\
&=\frac{\sqrt{8}}{5}\int_0^5{x\sqrt{1+\left( \frac{4}{5} \right) ^2}\mathrm{d}x}=\sqrt{82}.
\end{aligned}
\]
其中$\left( \cos \left( \widehat{\boldsymbol{t},x} \right) ,\sin \left( \widehat{\boldsymbol{t},y} \right) \right) $为积分路径曲线正切向的方向余弦.
没能成功的原因就是两类曲线积分转换时的损失,没有充分利用题目条件:当$y=x$时,$f$的两个偏导数都为0.上述证明中选取的积分路径与$y=x$这条直线关系不大.
\end{remark}
\begin{proof}
\begin{enumerate}
\item 注意到
\[
f(1,2) - f(0,0) = \int_{(0,0)}^{(1,2)} \left( \frac{\partial f}{\partial x} \mathrm{d}x + \frac{\partial f}{\partial y} \mathrm{d}y \right)
\]
这里积分路径待定. 于是由第一型曲线积分和第二型曲面积分的联系和\hyperref[theorem:Cauchy不等式(离散版)]{Cauchy不等式(离散版本)}得
\[
\begin{aligned}
|f(1,2)| &= \left| \int_{(0,0)}^{(1,2)} \left( \frac{\partial f}{\partial x} \cos \left( \widehat{\boldsymbol{t},x} \right)  + \frac{\partial f}{\partial y} \sin \left( \widehat{\boldsymbol{t},x} \right)  \right) \mathrm{d}s \right| \\
&\leqslant \int_{(0,0)}^{(1,2)} |\nabla f| \mathrm{d}s \leqslant \int_{(0,0)}^{(1,2)} 1 \mathrm{d}s
\end{aligned},
\]
其中$\left( \cos \left( \widehat{\boldsymbol{t},x} \right) ,\sin \left( \widehat{\boldsymbol{t},y} \right) \right) $为积分路径曲线正切向的方向余弦.为了得到这个方法最佳的估计, 我们取积分路径为连接 \( (0,0),(1,2) \) 的线段, 这恰好给出了 \( |f(1,2)| \leqslant \sqrt{5} \).

\item 先沿着 \( y = x, 0 \leqslant x \leqslant 4 \) 积分, 此时知积分 \( \int_{(0,0)}^{(4,4)} \left( \frac{\partial f}{\partial x} \mathrm{d}x + \frac{\partial f}{\partial y} \mathrm{d}y \right) \) 为 0. 于是
\[
\begin{aligned}
|f(5,4)| &= \left| \int_{(0,0)}^{(5,4)} \left( \frac{\partial f}{\partial x} \mathrm{d}x + \frac{\partial f}{\partial y} \mathrm{d}y \right) \right| = \left| \int_{(4,4)}^{(5,4)} \left( \frac{\partial f}{\partial x} \mathrm{d}x + \frac{\partial f}{\partial y} \mathrm{d}y \right) \right| \\
&= \left| \int_{4}^{5} \frac{\partial f}{\partial x} \mathrm{d}x \right| \leqslant 2 \int_{4}^{5} |x - 4| \mathrm{d}x = 1.
\end{aligned}
\]
\end{enumerate}
\end{proof}

\begin{example}
设$f(x,y) = x + y + xy$，求其在曲线$C: x^2 + y^2 + xy = 3$上的最大方向导数。
\end{example}
\begin{solution}
注意到$f(x,y)$在$C$上的方向导数最大值就是$f(x,y)$在$C$上梯度模的最大值。再由条件可得$f(x,y)$的梯度为
\begin{align*}
\nabla f = \left( \frac{\partial f}{\partial x}, \frac{\partial f}{\partial y} \right) = (1+y, 1+x).
\end{align*}
于是
\begin{align*}
|\nabla f|^2 = x^2 + y^2 + 2(x+y) + 2, \quad \forall (x,y) \in C.
\end{align*}
由$(x,y) \in C$知$x^2 + y^2 = 3 - xy$，令$u = x + y$, $v = xy$，则
\begin{align*}
x^2 + y^2 = 3 - xy \iff u^2 = 3 + v \iff v = u^2 - 3.
\end{align*}
进而
\begin{align*}
|\nabla f|^2 = 3 - xy + 2(x+y) + 2 = 5 - v + 2u = 9 - (u+1)^2.
\end{align*}
故
\begin{align*}
\max\limits_{(x,y) \in C} |\nabla f|^2 = 9,
\end{align*}
并且当$u = -1$时$|\nabla f|^2$取到最大值，即
\begin{align*}
u = x + y = -1,\ v = u^2 - 3 = -2 = xy \implies
\begin{cases}
x = -2, & y = 1 \\
\text{or} \\
x = 1, & y = -2
\end{cases}.
\end{align*}
因此$f$在$C$上的最大方向导数为
\begin{align*}
\max\limits_{(x,y) \in C} |\nabla f| = 3.
\end{align*}

\end{solution}


\subsection{多元函数的微分}

\begin{definition}
设$D \subseteq \mathbb{R}^n$为开集，$\bm{x}_0 \in D$，$f: D \to \mathbb{R}^m$。如果存在某个线性变换$A$（只依赖于$\bm{x}_0$），使得$\bm{x} \in U(\bm{x}_0) \subseteq D$时，有
\begin{align*}
f(\bm{x}) - f(\bm{x}_0) = A(\bm{x} - \bm{x}_0) + o\left( \|\bm{x} - \bm{x}_0\| \right)
\end{align*}
或
\begin{align*}
\lim\limits_{\bm{x} \to \bm{x}_0} \frac{f(\bm{x}) - f(\bm{x}_0) - A(\bm{x} - \bm{x}_0)}{\|\bm{x} - \bm{x}_0\|} = \bm{0},
\end{align*}
则称向量函数$f$在点$\bm{x}_0$\textbf{可微}（或\textbf{可导}）。若与上述线性变换$A$相联系的矩阵为$A_{m \times n}$，则称$A(\bm{x} - \bm{x}_0)$为$f$在点$\bm{x}_0$的\textbf{微分}，并称$A$为$f$在点$\bm{x}_0$的\textbf{导数}，记作$\mathrm{D}f(\bm{x}_0)$或$f'(\bm{x}_0)$。
\end{definition}

\begin{theorem}\label{theorem:多元函数可微的导数与偏导的关系}
设$D \subseteq \mathbb{R}^n$为开集，$\bm{x}_0 \in D$，$f: D \to \mathbb{R}^m$。若$f$在$\bm{x}_0$可微，则当$\bm{x} \in U(\bm{x}_0) \subseteq D$时，有
\begin{align*}
f(\bm{x}) - f(\bm{x}_0) = f'(x_0)(\bm{x} - \bm{x}_0) + o\left( \|\bm{x} - \bm{x}_0\| \right),
\end{align*}
其中
\begin{align*}
f'(\bm{x}_0) = \mathrm{D}f(\bm{x}_0) = \begin{pmatrix}
\frac{\partial f_1}{\partial x_1} & \cdots & \frac{\partial f_1}{\partial x_n} \\
\vdots & & \vdots \\
\frac{\partial f_m}{\partial x_1} & \cdots & \frac{\partial f_m}{\partial x_n}
\end{pmatrix}
.
\end{align*}
\end{theorem}

\begin{theorem}[多元函数带Peano余项的Taylor公式]\label{theorem:多元函数带Peano余项的Taylor公式}
设 $D \subseteq \mathbb{R}^n$ 是开集，$\bm{x}_0 \in D$，函数 $f: D \to \mathbb{R}$ 在 $\bm{x}_0$ 的某邻域内具有直到 $k$ 阶的连续偏导数,则对任意满足 $\bm{x}_0 + \bm{h} \in D$ 的 $\bm{h} = (h_1,\cdots,h_n)$，存在$\delta>0$,当 $\|\bm{h}\| < \delta$ 时,有
\[
f(\bm{x}_0 + \bm{h}) = \sum_{|\alpha|=0}^{k} \frac{1}{\alpha!} \, \partial^\alpha f(\bm{x}_0) \, \bm{h}^\alpha \;+\; o\bigl(\|\bm{h}\|^k\bigr),
\]
其中
\begin{enumerate}[(1)]
\item $\alpha = (\alpha_1,\cdots,\alpha_n)$ 是多重指标，$\alpha_i$ 为非负整数，$|\alpha| = \alpha_1 + \cdots + \alpha_n$；
\item $\alpha! = \alpha_1! \,\alpha_2! \cdots \alpha_n!$；
\item $\bm{h}^\alpha = h_1^{\alpha_1} h_2^{\alpha_2} \cdots h_n^{\alpha_n}$；
\item $\partial^\alpha f(\bm{x}_0) = \dfrac{\partial^{|\alpha|} f}{\partial x_1^{\alpha_1} \cdots \partial x_n^{\alpha_n}}(\bm{x}_0)$ 是混合偏导数在 $\bm{x}_0$ 处的值。
\end{enumerate}
等价地，用全微分形式可写作
\[
f(\bm{x}_0 + \bm{h}) = f(\bm{x}_0) + \sum_{j=1}^{k} \frac{1}{j!}\, \mathrm{d}^j f(\bm{x}_0)(\bm{h}) \;+\; o\bigl(\|\bm{h}\|^k\bigr),
\]
其中 $\mathrm{d}^j f(\bm{x}_0)$ 是 $j$ 阶全微分，作用在 $\bm{h}$ 上为
\[
\mathrm{d}^j f(\bm{x}_0)(\bm{h}) = \sum_{|\alpha|=j} \frac{j!}{\alpha!} \, \partial^\alpha f(\bm{x}_0) \, \bm{h}^\alpha.
\]

特别地，当$k=2$ 时，设$\boldsymbol{x}=\left( x_1,\cdots,x_n \right)^T ,\boldsymbol{x}_0=\left( x_1^0,\cdots,x_n^0 \right)^T $,$\bm{x_0} \in D \subseteq \mathbb{R}^n$，$f(\bm{x})$ 在 $\bm{x_0}$邻域内具有二阶连续偏导数，则存在$\delta>0$,当$\bm{h}\in \mathbb{R}^n$且$\|\bm{h}\|<\delta$时,有
\begin{align*}
f\left( \boldsymbol{x}_0+\boldsymbol{h} \right) =f\left( \boldsymbol{x}_0 \right) +\nabla f\left( \boldsymbol{x}_0 \right) \cdot \boldsymbol{h}+\frac{1}{2}\boldsymbol{h}^TH\left( \boldsymbol{x}_0 \right) \boldsymbol{h}+o\left( \left\| \boldsymbol{h} \right\| ^2 \right) .
\end{align*}
也即存在$\delta>0$,当$\bm{x}\in \mathbb{R}^n$且$\|\bm{x}\|< \delta$ 时有
\begin{align*}
f\left( \boldsymbol{x} \right) =f\left( \boldsymbol{x}_0 \right) +\nabla f\left( \boldsymbol{x}_0 \right) \cdot \left( \boldsymbol{x}-\boldsymbol{x}_0 \right) +\frac{1}{2}\left( \boldsymbol{x}-\boldsymbol{x}_0 \right) ^TH\left( \boldsymbol{x}_0 \right) \left( \boldsymbol{x}-\boldsymbol{x}_0 \right) +o\left( \left\| \left( \boldsymbol{x}-\boldsymbol{x}_0 \right) \right\| ^2 \right) ,
\end{align*}
其中$H\left( \boldsymbol{x}_0 \right)=\left( \begin{matrix}
f_{xx}\left( \boldsymbol{x}_0 \right)&		f_{xy}\left( \boldsymbol{x}_0 \right)\\
f_{yx}\left( \boldsymbol{x}_0 \right)&		f_{yy}\left( \boldsymbol{x}_0 \right)\\
\end{matrix} \right) $为$f$的二阶Hessian矩阵在$\bm{x_0}$处取值的矩阵.
\end{theorem}

\begin{example}
设$f: \mathbb{R}^2 \to \mathbb{R}^2$可微且恒有
\begin{align*}
\|f(x,y) - f(0,0)\| \geqslant \sqrt{x^2 + y^2}, \quad \forall (x,y) \in \mathbb{R}^2.
\end{align*}
证明：$f$在原点的Jacobi矩阵一定非退化，即行列式非零。
\end{example}
\begin{proof}
反证，若$|Jf(0)| = 0$，则存在$0 \ne X \in \mathbb{R}^2$使得$Jf(0)X = 0$. 由条件知
\begin{align*}
\|f(tX) - f(0)\| \geqslant \|tX\| = |t| \cdot \|X\|, \quad \forall t > 0.
\end{align*}
又由$f$在$\mathbb{R}^2$上可微和\refthe{theorem:多元函数可微的导数与偏导的关系}知
\begin{align*}
f(tX) = f(0) + Jf(0)tX + o(tX), \quad t \to 0^+.
\end{align*}
即对$\forall \varepsilon \in (0,1)$, $\exists \delta > 0$，使得当$t \in (0,\delta)$时有
\begin{align*}
\frac{\|f(tX) - f(0)\|}{\|tX\|} < \|Jf(0)\| + \varepsilon = \varepsilon < 1.
\end{align*}
这与前面的不等式矛盾！故结论成立.
\end{proof}

\begin{example}
设$f(x,y)$在平面$\mathbb{R}^2$上有二阶连续偏导数，$B((0,0),r)$是以$(0,0)$为心，$r$为半径的圆盘。证明：
\begin{align*}
\lim\limits_{r \to 0^+} \frac{1}{r^2} \left[ \frac{1}{\pi r^2} \iint_{B((0,0),r)} f(x,y) \, \mathrm{d}x\mathrm{d}y - f(0,0) \right] = \frac{1}{8} \Delta f(0,0),
\end{align*}
其中$\Delta f = f_{xx} + f_{yy}$。
\end{example}
\begin{proof}
由\hyperref[theorem:多元函数带Peano余项的Taylor公式]{多元函数带Peano余项的Taylor公式}知存在$\delta >0$，使得当$x^2+y^2<\delta$时，有
\begin{align}
f(x,y) &= f(0,0) + f_x(0,0)x + f_y(0,0)y + \frac{1}{2}\left[ f_{xx}(0,0)x^2 + 2f_{xy}(0,0)xy + f_{yy}(0,0)y^2 \right] + R(x,y), \label{eq9s7k2p}
\end{align}
并且$\lim\limits_{(x,y) \to (0,0)} \frac{R(x,y)}{x^2+y^2}=0.$从而对$\forall \varepsilon >0$，$\exists \delta_1 \in (0,\delta)$，使得当$x^2+y^2<\delta_1$时，有
\begin{align}
|R(x,y)| \leqslant \varepsilon (x^2+y^2). \label{eq4z8m1n}
\end{align}
对$\forall r \in (0,\delta)$，记$B_r \triangleq B((0,0),r)$，则由对称性知
\begin{gather*}
\iint_{B_r} x\mathrm{d}x\mathrm{d}y=\iint_{B_r} -x\mathrm{d}x\mathrm{d}y=\iint_{B_r} y\mathrm{d}x\mathrm{d}y=\iint_{B_r} -y\mathrm{d}x\mathrm{d}y=\iint_{B_r} xy\mathrm{d}x\mathrm{d}y=\iint_{B_r} -xy\mathrm{d}x\mathrm{d}y=0, \\
\iint_{B_r} x^2\mathrm{d}x\mathrm{d}y=\iint_{B_r} y^2\mathrm{d}x\mathrm{d}y=\frac{1}{2}\iint_{B_r} (x^2+y^2)\mathrm{d}x\mathrm{d}y=\frac{1}{2}\int_0^{2\pi}\mathrm{d}\theta \int_0^r \rho^3\mathrm{d}\rho=\frac{\pi r^4}{4}.
\end{gather*}
从而对$\eqref{eq9s7k2p}$式两边积分得
\begin{align}
\iint_{B_r} f(x,y)\mathrm{d}x\mathrm{d}y &= \iint_{B_r} \left[ f(0,0) + f_x(0,0)x + f_y(0,0)y + \frac{1}{2}\left[ f_{xx}(0,0)x^2 + 2f_{xy}(0,0)xy + f_{yy}(0,0)y^2 \right] + R(x,y) \right] \mathrm{d}x\mathrm{d}y \nonumber \\
&= f(0,0)\iint_{B_r}\mathrm{d}x\mathrm{d}y + \frac{f_{xx}(0,0)}{2}\iint_{B_r}x^2\mathrm{d}x\mathrm{d}y + \frac{f_{yy}(0,0)}{2}\iint_{B_r}y^2\mathrm{d}x\mathrm{d}y + \iint_{B_r}R(x,y)\mathrm{d}x\mathrm{d}y \nonumber \\
&= f(0,0)\pi r^2 + \frac{\pi r^4}{8}\left[ f_{xx}(0,0) + f_{yy}(0,0) \right] + \iint_{B_r}R(x,y)\mathrm{d}x\mathrm{d}y. \label{eq7d3s6f}
\end{align}
由$\eqref{eq4z8m1n}$式知当$r<\delta_1$时，有
\begin{align*}
\left| \frac{1}{r^4}\iint_{B_r} R(x,y)\mathrm{d}x\mathrm{d}y \right| \leqslant \frac{\varepsilon}{r^4}\iint_{B_r} (x^2+y^2)\mathrm{d}x\mathrm{d}y=\frac{\pi}{2}\varepsilon.
\end{align*}
故
\begin{align}
\lim\limits_{r\to 0^+} \frac{1}{r^4}\iint_{B_r} R(x,y)\mathrm{d}x\mathrm{d}y=0. \label{eq2p5g8h}
\end{align}
于是结合\eqref{eq9s7k2p}\eqref{eq7d3s6f}\eqref{eq2p5g8h}式可得
\begin{align*}
\lim\limits_{r\to 0^+} \frac{1}{r^2}\left[ \frac{1}{\pi r^2}\iint_{B_r} f(x,y)\mathrm{d}x\mathrm{d}y - f(0,0) \right] = \lim\limits_{r\to 0^+} \frac{1}{8}\left[ f_{xx}(0,0) + f_{yy}(0,0) \right] + \lim\limits_{r\to 0^+}\frac{1}{\pi r^4}\iint_{B_r}R(x,y)\mathrm{d}x\mathrm{d}y 
= \frac{1}{8}\Delta f(0,0).
\end{align*}
\end{proof}

\begin{example}
\(u(x, y, z)\) 是 \(R^3\) 上的连续函数, 在 \(u(x_0, y_0, z_0)\) 的某个邻域上具有二阶连续偏导. 定义
\begin{align*}
F(R) = \frac{1}{4\pi R^2} \iint_{\Sigma} u(x, y, z)\, \mathrm{d}S,
\end{align*}
其中\(\Sigma\) 是以 \((x_0, y_0, z_0)\) 为球心, 半径为 \(R\) 的球面.
\begin{enumerate}[(1)]
\item 证明 \(\lim\limits_{R \to 0} F(R) = u(x_0, y_0, z_0)\);
\item 证明 \(\lim\limits_{R \to 0} \frac{F(R) - u(x_0, y_0, z_0)}{R^2} = \frac{\Delta u(x_0, y_0, z_0)}{6}\), 其中 \(\Delta u = \frac{\partial^2 u}{\partial x^2} + \frac{\partial^2 u}{\partial y^2} + \frac{\partial^2 u}{\partial z^2}\).
\end{enumerate}
\end{example}
\begin{proof}
\begin{enumerate}[(1)]
\item 根据函数的连续性, 对于任意 \(\varepsilon > 0\), 存在 \(\delta > 0\), 当 \(\sqrt{(x-x_0)^2 + (y-y_0)^2 + (z-z_0)^2} < \delta\) 时, 有
\begin{align*}
|u(x, y, z) - u(x_0, y_0, z_0)| < \varepsilon.
\end{align*}
当 \(R < \delta\) 时, 球面$B(x_0,R)$上的所有点都满足上述条件,于是
\begin{align*}
\left|F(R) - u(x_0, y_0, z_0)\right| &= \left| \frac{1}{4\pi R^2} \iint_{\Sigma} \left[u(x, y, z) - u(x_0, y_0, z_0)\right]\, \mathrm{d}S \right| \\
&\leqslant \frac{1}{4\pi R^2} \iint_{\Sigma} \left|u(x, y, z) - u(x_0, y_0, z_0)\right|\, \mathrm{d}S \\
&< \frac{1}{4\pi R^2} \cdot \varepsilon \cdot 4\pi R^2 = \varepsilon .
\end{align*}
因此, 由极限定义得\(\lim\limits_{R \to 0} F(R) = u(x_0, y_0, z_0)\).

\item 设 \(u_0 = u(x_0, y_0, z_0)\). 根据\hyperref[theorem:多元函数带Peano余项的Taylor公式]{多元函数的Taylor公式}得
\begin{align*}
u(x, y, z) = u_0 + \nabla u_0 \cdot \bm{r} + \frac{1}{2} \bm{r}^T H u_0 \bm{r} + o(R^2),
\end{align*}
其中 \(\bm{r} = (x, y, z)\), \(|\bm{r}| = R\),且\(Hu_0\) 是$u$在$(x_0,y_0,z_0)$处取值的Hessian矩阵.
在球面$\Sigma$上对 \(u(x, y, z)\) 进行积分得
\begin{align}
F(R)=\iint_{\Sigma} u(x, y, z)\, \mathrm{d}S = \iint_{\Sigma} u_0\, \mathrm{d}S + \iint_{\Sigma} (\nabla u_0 \cdot \bm{r})\, \mathrm{d}S + \iint_{\Sigma} \frac{1}{2} \bm{r}^T H u_0 \bm{r}\, \mathrm{d}S + \iint_{\Sigma} o(R^2)\, \mathrm{d}S .\label{eq::jfijeioWQQQ}
\end{align}
其中
\begin{align*}
\iint_{\Sigma} u_0\, \mathrm{d}S = u_0 \cdot 4\pi R^2.
\end{align*}
由于球面关于原点对称,故
\begin{align}\label{eq::fjisfjijJFEOEO}
\iint_{\Sigma}{x\,\mathrm{d}S}=\iint_{\Sigma}{y\,\mathrm{d}S}=\iint_{\Sigma}{z\,\mathrm{d}S}=\iint_{\Sigma}{xy\,\mathrm{d}S}=\iint_{\Sigma}{xz\,\mathrm{d}S}=\iint_{\Sigma}{yz\,\mathrm{d}S}=\iint_{\Sigma}{xyz\,\mathrm{d}S}=0.
\end{align}
从而
\begin{align*}
\iint_{\Sigma} (\nabla u_0 \cdot \bm{r})\, \mathrm{d}S = 0.
\end{align*}
再利用球坐标积分对称性知
\begin{align}\label{eq::fjisfjijJFEOEO-1}
\iint_{\Sigma} x^2\, \mathrm{d}S = \iint_{\Sigma} y^2\, \mathrm{d}S = \iint_{\Sigma} z^2\, \mathrm{d}S = \frac{1}{3} \iint_{\Sigma} (x^2+y^2+z^2)\, \mathrm{d}S = \frac{1}{3} \iint_{\Sigma} R^2\, \mathrm{d}S = \frac{4\pi R^4}{3}.
\end{align}
再由\eqref{eq::fjisfjijJFEOEO}\eqref{eq::fjisfjijJFEOEO-1}式可得
\begin{align*}
\iint_{\Sigma} \bm{r}^T H u_0 \bm{r}\, \mathrm{d}S = \sum\limits_{i=1}^3 u_{ii} \cdot \frac{4\pi R^4}{3} = \frac{4\pi R^4}{3} \Delta u_0.
\end{align*}
且
\begin{align*}
\iint_{\Sigma} o(R^2)\, \mathrm{d}S = o(R^2) \cdot 4\pi R^2 = o(R^4).
\end{align*}
将上述各项代入\eqref{eq::jfijeioWQQQ}式得
\begin{align*}
F(R) = \frac{1}{4\pi R^2} \left( 4\pi R^2 u_0 + 0 + \frac{1}{2} \cdot \frac{4\pi R^4}{3} \Delta u_0 + o(R^4) \right).
\end{align*}
化简得
\begin{align*}
F(R) = u_0 + \frac{R^2}{6} \Delta u_0 + o(R^2).
\end{align*}
因此
\begin{align*}
\frac{F(R) - u_0}{R^2} = \frac{\Delta u_0}{6} + \frac{o(R^2)}{R^2}.
\end{align*}
令\(R \to 0\)得
\begin{align*}
\lim\limits_{R \to 0} \frac{F(R) - u(x_0, y_0, z_0)}{R^2} = \frac{\Delta u(x_0, y_0, z_0)}{6}.
\end{align*}
\end{enumerate}
\end{proof}

\begin{theorem}[拟微分平均值定理]
设凸区域 $D\subset \mathbf{R}^{n}$，$f$ 在 $D$ 上可微，$a,b\in D$，则 $\exists \theta \in (0,1)$ 使
\begin{align*}
|f(b) - f(a)| \leqslant \| Jf(a + \theta (b - a))\| \cdot |b - a|,
\end{align*}
这里 $\| \cdot \|$ 表示矩阵的模，即
\begin{align*}
\| A\| = \max_{|x| = 1}|Ax|.
\end{align*}
\end{theorem}
\begin{proof}
用$\langle \cdot, \cdot \rangle$ 表示 $\mathbf{R}^m$ 上的标准Euclid内积，即对任意 $x=(x_1,\cdots,x_m),\, y=(y_1,\cdots,y_m)\in \mathbf{R}^m$，有
\begin{align*}
\langle x,y\rangle = \sum_{i=1}^m x_i y_i.
\end{align*}
设
\begin{align*}
\phi (t) = \langle f(a + t(b - a)) - f(a),\, f(b) - f(a)\rangle ,\quad t\in [0,1].
\end{align*}
由\hyperref[Functional Analysis-proposition:内积的连续性]{内积的连续性}知,极限和内积可交换,进而求导也和内积可交换.于是对 $\phi$ 求导，再由链式法则得
\begin{align*}
\phi'(t) = \left\langle \frac{\mathrm{d}}{\mathrm{d}t} f(a+t(b-a)),\, f(b)-f(a) \right\rangle = \langle Jf(a+t(b-a))(b-a), \,f(b)-f(a)\rangle ,
\end{align*}
其中$Jf$ 表示 $f$ 的Jacobi矩阵,即
\begin{align*}
Jf\triangleq \left(\frac{\partial f_i}{\partial x_j}\right)_{m\times n}.
\end{align*}
由一元函数的拉格朗日中值定理，存在 $\theta \in (0,1)$，使得
\begin{align*}
|f(b) - f(a)|^2 =\phi(1) - \phi(0) = \phi'(\theta)= \langle Jf(a+\theta(b-a))(b-a), \,f(b)-f(a)\rangle .
\end{align*}
再由Cauchy-Schwarz不等式得
\begin{align*}
|f(b) - f(a)|^2 \leqslant |Jf(a + \theta (b - a))(b - a)|\, |f(b) - f(a)|.
\end{align*}
再由\hyperref[Functional Analysis-proposition:有界线性算子的基本不等式]{这里矩阵模的定义}，有
\begin{align*}
|Ax| \leqslant \|A\|\, |x|,\quad \forall x\in D.
\end{align*}
因此
\begin{align*}
|Jf(a + \theta (b - a))(b - a)| \leqslant \|Jf(a + \theta (b - a))\|\, |b - a|.
\end{align*}
代入得
\begin{align*}
|f(b) - f(a)|^2 \leqslant \|Jf(a + \theta (b - a))\|\, |b - a|\, |f(b) - f(a)|.
\end{align*}
若 $f(b)=f(a)$，则结论显然。否则 $|f(b)-f(a)|>0$，两边除以它，得到
\begin{align*}
|f(b) - f(a)| \leqslant \|Jf(a + \theta (b - a))\|\, |b - a|.
\end{align*}
\end{proof}



\subsection{多元函数极值原理}

\begin{theorem}[多元函数极值原理]\label{theorem:多元函数极值原理}
设 \( f \) 是某个区域 \( V \subset \mathbb{R}^n \) 的二阶连续可微函数, 我们定义其 Hess(黑塞) 矩阵为
\[
Hf = \left( \frac{\partial^2 f}{\partial x_i \partial x_j} \right)_{1 \leqslant i, j \leqslant n},
\]
则对 \( \mathbf{x}_0 \in V \) 满足 \( \frac{\partial f}{\partial x_i}(\mathbf{x}_0) = 0, i = 1, 2, \cdots, n \) 有
\begin{enumerate}[(1)]
\item \( Hf(\mathbf{x}_0) \) 是正定的, 则 \( \mathbf{x}_0 \) 是 \( f \) 严格极小值点;

\item \( Hf(\mathbf{x}_0) \) 是负定的, 则 \( \mathbf{x}_0 \) 是 \( f \) 严格极大值点;

\item \( Hf(\mathbf{x}_0) \) 是不定的(既不是正定,也不是负定), 则 \( \mathbf{x}_0 \) 不是 \( f \) 极值点;

\item 若 \( \mathbf{x}_0 \) 是 \( f \) 极小值点, 则 \( Hf(\mathbf{x}_0) \) 是半正定的;

\item 若 \( \mathbf{x}_0 \) 是 \( f \) 极大值点, 则 \( Hf(\mathbf{x}_0) \) 是半负定的.
\end{enumerate}
\end{theorem}
\begin{proof}
令 \( \bm{h} = \bm{x} - \bm{x}_0 \). 由于 \( f \) 二阶连续可微, 由\hyperref[theorem:多元函数带Peano余项的Taylor公式]{多元函数带Peano余项的Taylor公式}可得
\begin{align*}
f(\bm{x}_0 + \bm{h}) = f(\bm{x}_0) + \nabla f(\bm{x}_0) \cdot \bm{h} + \frac{1}{2} \bm{h}^T Hf(\bm{x}_0) \bm{h} + o(\|\bm{h}\|^2). 
\end{align*}
又由题设知 \( \nabla f(\bm{x}_0) = 0 \), 故上式可化为
\begin{align}
f(\bm{x}_0 + \bm{h}) - f(\bm{x}_0) = \frac{1}{2} \bm{h}^T Hf(\bm{x}_0) \bm{h} + o(\|\bm{h}\|^2). \label{eq:taylor_simplified}
\end{align}
记 \( H = Hf(\bm{x}_0) \).
\begin{enumerate}[(1)]
\item 若 \( H \) 正定, 因为单位球面 \( S^{n-1} = \{ \bm{\xi} \in \mathbb{R}^n \mid \|\bm{\xi}\| = 1 \} \) 是紧集,$\bm{x}\mapsto \bm{x}^TH\bm{x}$是连续函数,所以由\hyperref[Functional Analysis-proposition:紧集上的连续函数必有最值]{紧集上的连续函数必有最值}知,存在常数
\begin{align*}
m = \min_{\|\bm{\xi}\| = 1} \bm{\xi}^T H \bm{\xi} > 0.
\end{align*}
于是对任意 \( \bm{h} \in \mathbb{R}^n \), 均有
\begin{align}
\bm{h}^T H \bm{h} \geqslant m \|\bm{h}\|^2. \label{eq:positive_definite_lower_bound}
\end{align}
对于无穷小量 \( o(\|\bm{h}\|^2) \), 存在 \( \delta > 0 \), 使得当 \( 0 < \|\bm{h}\| < \delta \) 时, 有 \( |o(\|\bm{h}\|^2)| < \frac{m}{4} \|\bm{h}\|^2 \). 结合\eqref{eq:taylor_simplified}\eqref{eq:positive_definite_lower_bound}式, 可得
\begin{align*}
f(\bm{x}_0 + \bm{h}) - f(\bm{x}_0)  \geqslant \frac{1}{2} m \|\bm{h}\|^2 - |o(\|\bm{h}\|^2)| 
> \frac{1}{2} m \|\bm{h}\|^2 - \frac{1}{4} m \|\bm{h}\|^2 = \frac{1}{4} m \|\bm{h}\|^2 > 0.
\end{align*}
因此\( \bm{x}_0 \) 是 \( f \) 的严格极小值点.

\item 若 \( H \) 负定, 设
\begin{align*}
M = \max_{\|\bm{\xi}\| = 1} \bm{\xi}^T H \bm{\xi} < 0.
\end{align*}
则对任意 \( \bm{h} \), 有
\begin{align}
\bm{h}^T H \bm{h} \leqslant M \|\bm{h}\|^2 < 0. \label{eq:negative_definite_upper_bound}
\end{align}
同样存在 \( \delta > 0 \) 使得当 \( 0 < \|\bm{h}\| < \delta \) 时, \( |o(\|\bm{h}\|^2)| < -\frac{1}{4} M \|\bm{h}\|^2 \). 于是结合 \eqref{eq:taylor_simplified}\eqref{eq:negative_definite_upper_bound}式, 可得
\begin{align*}
f(\bm{x}_0 + \bm{h}) - f(\bm{x}_0) \leqslant \frac{1}{2} M \|\bm{h}\|^2 + |o(\|\bm{h}\|^2)| 
< \frac{1}{2} M \|\bm{h}\|^2 - \frac{1}{4} M \|\bm{h}\|^2 
= \frac{1}{4} M \|\bm{h}\|^2 < 0.
\end{align*}
因此\( \bm{x}_0 \) 是 \( f \) 的严格极大值点.
\item 若 \( H \) 不定, 则存在 \( \bm{u}, \bm{v} \in \mathbb{R}^n \) 使得 \( \bm{u}^T H \bm{u} > 0 \) 且 \( \bm{v}^T H \bm{v} < 0 \). 令 \( \bm{h} = t \bm{u} \), 当 \( t \to 0 \) 且 \( t \neq 0 \) 时, 由 \eqref{eq:taylor_simplified}式可知
\begin{align*}
f(\bm{x}_0 + t \bm{u}) - f(\bm{x}_0) = \frac{1}{2} t^2 \bm{u}^T H \bm{u} + o(t^2) > 0.
\end{align*}
令 \( \bm{h} = t \bm{v} \), 当 \( t \to 0 \) 且 \( t \neq 0 \) 时,同样由 \eqref{eq:taylor_simplified}式可知
\begin{align*}
f(\bm{x}_0 + t \bm{v}) - f(\bm{x}_0) = \frac{1}{2} t^2 \bm{v}^T H \bm{v} + o(t^2) < 0.
\end{align*}
即在 \( \bm{x}_0 \) 的任意去心邻域内均存在两点, 其函数值分别大于和小于 \( f(\bm{x}_0) \), 故 \( \bm{x}_0 \) 不是 \( f \) 的极值点.
\item 反证法. 假设 \( \bm{x}_0 \) 是 \( f \) 的极小值点, 但 \( H \) 不是半正定的. 则存在 \( \bm{w} \in \mathbb{R}^n \) 使得 \( \bm{w}^T H \bm{w} < 0 \). 取 \( \bm{h} = t \bm{w} \), 令 \( t \to 0 \), 代入 \eqref{eq:taylor_simplified}式得
\begin{align*}
f(\bm{x}_0 + t \bm{w}) - f(\bm{x}_0) = \frac{1}{2} t^2 \bm{w}^T H \bm{w} + o(t^2) < 0,
\end{align*}
这与 \( \bm{x}_0 \) 是极小值点矛盾! 因此 \( H \) 必为半正定.

\item 反证法. 假设 \( \bm{x}_0 \) 是 \( f \) 的极大值点, 但 \( H \) 不是半负定的. 则存在 \( \bm{w} \in \mathbb{R}^n \) 使得 \( \bm{w}^T H \bm{w} > 0 \). 取 \( \bm{h} = t \bm{w} \), 令 \( t \to 0 \), 代入 \eqref{eq:taylor_simplified}式得
\begin{align*}
f(\bm{x}_0 + t \bm{w}) - f(\bm{x}_0) = \frac{1}{2} t^2 \bm{w}^T H \bm{w} + o(t^2) > 0,
\end{align*}
这与 \( \bm{x}_0 \) 是极大值点矛盾! 因此 \( H \) 必为半负定.
\end{enumerate}
\end{proof}

\begin{example}
设$f\in C^2(\mathbb{R}^2)$满足
\begin{align*}
\lim\limits_{R\to+\infty}\sup\limits_{x^2+y^2=R^2}|f(x,y)|=0,\ f(0,0)=0;
\end{align*}
\begin{align*}
x^2f_{xx}''+y^2f_{yy}''+e^xf_x'+y^3f_y'=(x^4+y^2)f.
\end{align*}
证明: $f\equiv0$.
\end{example}
\begin{proof}
反证.不妨设$f(x_0,y_0)>0$,其中$(x_0,y_0)\in\mathbb{R}^2$.由题设知,$\exists R_0>0$,使得
\begin{align*}
\sup\limits_{x^2+y^2=R^2}|f(x,y)|<\frac{f(x_0,y_0)}{2},\quad\forall R\geqslant R_0\Longrightarrow \sup\limits_{(x,y)\in\mathbb{R}^2\setminus B(0,R_0)}|f(x,y)|\leqslant\frac{f(x_0,y_0)}{2}.
\end{align*}
因为$B(0,R_0)$是有界闭集,所以$\exists (a,b)\in B(0,R_0)$,使得
\begin{align*}
f(a,b)=\max\limits_{(x,y)\in B(0,R_0)}f(x,y)\geqslant f(x_0,y_0)>\sup\limits_{(x,y)\in\mathbb{R}^2\setminus B(0,R_0)}|f(x,y)|.
\end{align*}
故
\begin{align*}
\max\limits_{(x,y)\in\mathbb{R}^2}f(x,y)=f(a,b)>0.
\end{align*}
因此$f(a,b)$也是$f(x,b),f(a,y)$在$\mathbb{R}$上的最大值,于是
\begin{align*}
f_x'(a,b)=\left.\frac{\mathrm{d}f(x,b)}{\mathrm{d}x}\right|_{x=a}=0,\quad f_y'(a,b)=\left.\frac{\mathrm{d}f(a,y)}{\mathrm{d}y}\right|_{y=b}=0.
\end{align*}
若$f_{xx}''(a,b)=\left.\frac{\mathrm{d}^2f(x,b)}{\mathrm{d}x^2}\right|_{x=a}>0$,则由\refthe{theorem:多元函数极值原理}知$f(x,b)$在$x=a$处取得极小值,这与$(a,b)$也是$f(x,b)$在$\mathbb{R}$上的最大值矛盾!故$f_{xx}''(a,b)\leqslant 0$.同理可知$f_{yy}''(a,b)\leqslant 0$.从而再由题设可知
\begin{align*}
0\geqslant a^2f_{xx}''(a,b)+b^2f_{yy}''(a,b)=(a^4+b^2)f(a,b)\geqslant 0.
\end{align*}
故$a=b=0$,即$f(0,0)>0$,这与$f(0,0)=0$矛盾!
\end{proof}

\begin{example}
已知$u(x,y)$在$D=\{(x,y)\mid x^2+y^2\leqslant1\}$上连续，且在$x^2+y^2<1$满足
\begin{align*}
\frac{\partial^2 u}{\partial x^2}+\frac{\partial^2 u}{\partial y^2}=u.
\end{align*}
\begin{enumerate}[(1)]
\item 证明: 若在$x^2+y^2=1$上有$u(x,y)\geqslant0$, 则在$x^2+y^2\leqslant1$上也有$u(x,y)\geqslant0$;
\item 证明: 若在$x^2+y^2=1$上$u(x,y)>0$, 则在$x^2+y^2\leqslant1$上也有$u(x,y)>0$.
\end{enumerate}
\end{example}
\begin{proof}
\begin{enumerate}[(1)]
\item 反证法. 由于$u$在$D$上连续, 所以$u$在$D$上存在最小值, 设最小值点为$(a,b)$. 若存在$(x_0,y_0)\in D$, 使得$u(x_0,y_0)<0$, 则$u(a,b)\leqslant u(x_0,y_0)<0$, 进而$(a,b)\in D^\circ$, 故
\begin{align*}
\frac{\partial u}{\partial x}(a,b)=\frac{\partial u}{\partial y}(a,b)=0.
\end{align*}
若$\frac{\partial ^2u}{\partial x^2}\left( a,b \right) =\left.\frac{\mathrm{d}^2u\left( x,b \right)}{\mathrm{d}x^2}\right|_{x=a}<0$, 那么根据\refthe{theorem:多元函数极值原理}可知$u(x,b)$在$(a,b)$处取得极大值, 这与$u(x,y)$在$(a,b)$处取得极小值矛盾!所以$\frac{\partial^2 u}{\partial x^2}(a,b)\geqslant0$, 同理也有$\frac{\partial^2 u}{\partial y^2}(a,b)\geqslant0$. 从而结合已知有
\begin{align*}
u(a,b)=\frac{\partial^2 u}{\partial x^2}(a,b)+\frac{\partial^2 u}{\partial y^2}(a,b)\geqslant0,
\end{align*}
这与$u(a,b)<0$矛盾!所以$x^2+y^2\leqslant1$上也有$u(x,y)\geqslant0$.

\item 由于$u$在有界闭集$\partial D: x^2+y^2=1$上恒正连续, 所以存在最小值, 设最小值为$m>0$. 而函数$e^x+e^y$也在$\partial D$上恒正连续, 所以存在最大值, 设最大值为$M>0$, 取正数$c=\frac{m}{M}$, 那么$m-cM=0$. 现构造函数$v(x,y)=u(x,y)-c(e^x+e^y)$, 显然在$v$也在$D$上连续, 且在$x^2+y^2<1$满足
\begin{align*}
\frac{\partial^2 v}{\partial x^2}+\frac{\partial^2 v}{\partial y^2}
=\frac{\partial^2 u}{\partial x^2}+\frac{\partial^2 u}{\partial y^2}-c(e^x+e^y)
=u-c(e^x+e^y)=v.
\end{align*}
同时在$x^2+y^2=1$上还有$v(x,y)\geqslant m-cM=0$, 由(1)可知在$x^2+y^2\leqslant1$上也有$v(x,y)\geqslant0$, 即
\begin{align*}
u(x,y)\geqslant c(e^x+e^y)>0.
\end{align*}
\end{enumerate}
\end{proof}

\begin{example}
设$f(x,y)=ax^2+2bxy+cy^2+dx+ey$, $x,y\in[0,1]^2\triangleq D$. 若$f(x,y)\geqslant0$, $(x,y)\in\partial D$, 证明: $f(x,y)\geqslant0$, $(x,y)\in D$.
\end{example}
\begin{proof}
由条件知
\begin{gather*}
f\left( x,0 \right) =x\left( ax+d \right) \geqslant 0,\forall x\in \left[ 0,1 \right] \Longrightarrow ax+d\geqslant 0,\forall x\in \left( 0,1 \right] .
\\
f\left( 0,y \right) =y\left( cy+e \right) \geqslant 0,\forall y\in \left[ 0,1 \right] \Longrightarrow cy+e\geqslant 0,\forall y\in \left( 0,1 \right] .
\end{gather*}
令$x,y\rightarrow 0^+$得$d,e\geqslant 0$. 显然$f\in C(D)$,所以$f$在$D$上存在最小值.若$f$的最小值点在$D$的边界达到, 则结论已证毕! 下设$(x_0,y_0)\in D^\circ$为$f$的最小值点, 则$x=x_0$和$y=y_0$分别是$f(x,y_0)$和$f(x_0,y)$的极小值点.从而
\begin{align*}
\frac{\partial f}{\partial x}\left( x_0,y_0 \right) =\left.\frac{\mathrm{d}f\left( x,y_0 \right)}{\mathrm{d}x}\right|_{x=x_0}=2ax_0+2by_0+d=0,\quad \frac{\partial f}{\partial y}\left( x_0,y_0 \right) =\left.\frac{\mathrm{d}f\left( x_0,y \right)}{\mathrm{d}y}\right|_{y=y_0}=2bx_0+2cy_0+e=0。
\end{align*}
解得$d=-2ax_0-2by_0,e=-2bx_0-2cy_0$.于是
\begin{align*}
f(x_0,y_0)=ax_0^2+2bx_0y_0+cy_0^2+dx_0+ey_0=-ax_0^2-cy_0^2-2bx_0y_0=dx_0+ey_0-f(x_0,y_0),
\end{align*}
即$f(x_0,y_0)=\frac{1}{2}(dx_0+ey_0)\geqslant0$. 故
\begin{align*}
f(x,y)\geqslant f(x_0,y_0)\geqslant 0,\quad \forall (x,y)\in D.
\end{align*}
\end{proof}














\end{document}